\section{Continuity across jobs, interfaces, and ownership}
\label{exec17:integration-continuity}

The local service already demonstrates several useful engineering ideas:
frozen requests, context-bound results, bounded workers, durable report
checkpoints, and matching responses across selected local interfaces. The next
proposal task is to state exactly what those observations warrant and what they
do not. A banking analyst may submit a portfolio-scale summary, inspect a
conditional treatment effect through a library, invoke the same request from a
command line, and review it in an agent-mediated workflow. Correctness requires
that a restart keeps the original scientific question, that each interface
preserves the meaning of results and refusals, that protocol and data schemas
are versioned independently, and that an institutional owner—not the author—
accepts a precisely identified release. These are related requirements; none
implies that another is already qualified.

The examples below concern a servicing report that summarizes canonical
responses for a frozen set of eligible accounts. “Actual” refers only to the
saved local executions on generated inputs. “Static” refers to source
inspection. “Authored reference” denotes a reproducible mathematical or finite
decision-table calculation. “Proposed” denotes a design awaiting its stated
acceptance tests. None is a bank result.\citep{exec16_parity_observed,exec14_failure_observations}

\subsection{A durable report is not an arbitrary resumable estimator}
\label{exec17:continuity-jobs}

The report worker has a concrete recovery invariant. Let the submitted ordered
artifact identities be $(a_1,\ldots,a_m)$ and let $C_k$ be the ordered
identities already summarized. A valid checkpoint has
\[
 C_k=(a_1,\ldots,a_k),\qquad
 h_k=H(J,C_k,S_k),
\]
where $J$ binds the original request, execution identity, artifact hashes,
and declared budget, and $S_k$ is the corresponding ordered summary list.
Resume checks that the identities form this prefix, that there is one summary
per completed identity, and that the checkpoint digest matches. It refuses a
changed execution identity or checkpoint, caps starts at three, and computes
remaining wall time from the original creation time. Those checks protect a
declared report reconstruction. The worker summarizes response declarations
under a schema; it explicitly does not validate the causal claims inside them.
\citep{exec14_failure_implementation}

The saved fault trial killed a report worker after a checkpoint and later
observed a completed report with one entry per frozen input. A distinct
instrumented trial interrupted storage after content publication but before
artifact metadata; same-content retry completed the descriptor, while the
content-only orphan remained unlisted in the meantime. Another trial changed
checkpoint content and observed a resume refusal. These are actual local
observations with generated requests and controlled interruption hooks. The
hooks establish behavior at named write boundaries; they are not naturally
observed host or storage failures. The records cover one-host report
reconstruction, not arbitrary estimator continuation, multi-host fencing, or
exactly-once effects at an external bank service.
\citep{exec14_failure_observations,exec13_ops_faults}

There are at least three persistence transitions: a per-input checkpoint, the
content and descriptor writes for a content-addressed artifact, and the job's
terminal acknowledgement. A lock protects ordinary in-process transitions,
but a process crash releases it. If result bytes and metadata are durable and
the process fails before the terminal state is written, a later cancellation
may make the job terminally cancelled while retaining a report artifact that
was already published. This does not violate terminal immutability: the
terminal job state and an independently addressable immutable artifact are
different objects. The client must read the job handle and artifact reference
as separate authorities. Cancellation does not roll back durable bytes, sent
messages, or other effects. The current causalbank job operation applies only
to schema-summary reports; it is not MCP request cancellation or an arbitrary
scientific worker API. The legacy protocol itself distinguishes request and
task cancellation and requires parties to handle completion races.
\citep{v12_mcp_legacy_cancellation,exec14_failure_implementation}

An exact finite reference enumerates all sixteen ordered selections without
repetition of three checkpoint identities at lengths zero through three.
Exactly four are valid prefixes, including the empty prefix. It also separates
five lock-respecting cancellation placements from the crash window between
publication and terminal acknowledgement. In that window the proposed recovery
checks the exact input, report, and job-state identities, preserves the
immutable artifact, and lets an already-recorded terminal state win. If a
descriptor points to missing or mismatched bytes, the proposal refuses
automatic repair and retains the evidence for operator reconciliation. These
counts are combinatorial checks of the stated finite model, not failure rates
or measurements of the implementation.
\citep{exec17_continuity_reference}

The deadline creates a separate static concern. Current resume logic computes
\[
 r_{\rm wall}=t_{\rm created}+B-t_{\rm now},
\]
where $B$ is the original budget and both $t$ values are wall-clock readings.
With $t_{\rm created}=1000$, $B=60$, and fifty seconds of real elapsed time,
the formula returns ten seconds. If eighty seconds have actually elapsed while
the clock has moved backward ninety seconds, the same formula returns seventy
seconds, even though the real elapsed budget has expired. This is a
mathematical counterexample to the formula under a wall-clock adjustment, not
an observed clock failure. A monotonic clock is the proper within-process
measure; a reboot-safe cross-process deadline needs an explicit trusted time
source or a conservative failure policy, because a monotonic epoch is not
portable across restarts. The accepted design must choose and test that
policy rather than silently promise both properties.
\citep{exec14_failure_implementation,exec17_continuity_reference}

The remaining local acceptance matrix should freeze tests for every worker
boundary: before and after each checkpoint; content before metadata; metadata
before terminal acknowledgement; cancellation before work, during work,
after publication, and after completion; repeated resume requests; stale
leases; process death; malformed or truncated state; changed code, dependency,
schema, artifact, or launch identity; deadline expiry; backward and forward
clock changes; and duplicate external-effect attempts. Each trial must compare
the full request identity, ordered prefix, artifact bytes and descriptor,
terminal state, attempt count, and emitted effects. Expected refusals belong in
the denominator. A fault harness may force the process boundary, but the
target platform owner must later run the frozen protocol on its declared
storage, scheduler, hardware, clock, and recovery topology. Current external
input contracts leave the platform operator and recovery lead unassigned.
\citep{exec17_continuity_contract}

\subsection{One semantic operation domain across three surfaces}
\label{exec17:shared-operation}

Let $\mathcal Q_L$, $\mathcal Q_C$, and $\mathcal Q_M$ denote requests admitted
by the library, CLI, and MCP tool map. The defensible parity target is the
negotiated intersection
\[
 \mathcal Q_* = \mathcal Q_L\cap\mathcal Q_C\cap\mathcal Q_M,
 \qquad
 L(q)=C(q)=M_{\rm structured}(q)=M_{\rm text}(q),\quad q\in\mathcal Q_*.
\]
Equality means the same canonical response after transport decoding: request
and context hashes, operation, invocation status, scientific disposition,
units, estimate or identified set, topology, uncertainty meaning, evidence,
qualification, and typed error. It does not mean identical JSON-RPC framing,
CLI exit code, or tool-level error wrapper. Those belong to different layers.
The MCP server must expose only operations in its actual mapping; the generic
runtime may return a canonical unsupported result for an operation the tool
map does not advertise.\citep{exec16_parity_design,exec16_parity_observed}

The actual saved parity suite ran thirteen generated scientific requests and
three lifecycle comparisons through the local library, CLI, and MCP stdio
client/server. Complete canonical results agreed in the tested common domain;
all nine advertised operations were exercised. Separate adverse tests observed
an unadvertised tool, schema rejection, and a valid request whose corrupted
stored replay was mislabeled as invalid caller input by broad adapter error
handling. The latter is a concrete error-origin defect, not a causal refusal.
The suite establishes these selected cases for the recorded Python and SDK
versions. It does not establish full parity over every schema branch, a newer
SDK, another language, HTTP transport, or a bank consumer.
\citep{exec16_parity_observed}

A shared-operation design should have one versioned registry that binds the
canonical request schema, response validator, method identifier, common
capability, and typed failure origin. The library runtime owns dispatch and
scientific context checks. The CLI owns bounded input parsing, command-name
agreement, JSON serialization, and documented exit mapping. MCP owns the
negotiated tool list, wire protocol, and error wrapper. Each adapter must call
the same canonical runtime and validate the same response object before
encoding it. A protocol-level “unknown tool” remains a capability error;
canonical unsupported analysis remains a scientific disposition; malformed
input remains an input error; corrupt replay remains stored-state failure;
and worker death remains worker failure. A successful transport is not a
scientific success. This design preserves semantic ownership while preventing
each adapter from growing an independent operation implementation.
\citep{exec16_parity_design,exec14_failure_design}

For an account-value request, changing the treatment sign, the thirty-day
horizon, USD units, or independent-account sampling law changes the scientific
target even if the JSON schema version remains unchanged. The context identity
must therefore change, and a previous moment or audit artifact bound to the
old context must be rejected. Conversely, changing a CLI wrapper without
changing the target should require adapter conformance evidence without
misrepresenting it as a new causal estimand. A reproducible authored example
uses a stipulated effect set $[-2,6]$ USD per account and a cost of USD $3$.
For $1000$ accounts the net-value set is
\[
 1000\bigl([-2,6]-3\bigr)=[-5000,3000]\text{ USD}.
\]
Converting both effect and cost to thousand-dollar units and then converting
the result back leaves that set unchanged. Converting only one quantity does
not. This exact interval arithmetic demonstrates why response equality must
preserve units and sign; it is not an estimate or a profit claim.
\citep{exec17_continuity_reference}

The repair sequence is concrete: define the registry and common domain;
separate caller, protocol, stored-state, artifact, worker, and transport
failures; freeze request/error pairs before running; test full-object equality
on every adapter branch; include changed context and artifact mismatches;
then add actual maintained consumers one at a time. A consumer's decoder must
preserve null endpoints, disconnected sets, large integers, explicit refusal,
and uncertainty labels. Current inputs identify consumer repositories and
platform/security owners but provide neither their actual decoder contracts
nor signed workflow acceptance. Language adapters can be specified now;
their external acceptance cannot be authored locally.
\citep{exec16_parity_design,exec17_continuity_contract}

\subsection{Version the scientific contract separately from its transport}
\label{exec17:versions}

“Version” refers to several identities. The proposed compatibility manifest
records these independently:
\begin{center}
\begin{tabular}{p{0.19\linewidth}p{0.28\linewidth}p{0.43\linewidth}}
\hline
Identity & Example & Change that requires review\\
\hline
Scientific context & Target, sign, horizon, units, assumptions, clocks, sampling & Changes estimand or interpretation; invalidate context-bound evidence\\
Request/response schema & \texttt{request/1.0.0} & Changes fields, constraints, or representation; migrate and run schema tests\\
Operation contract & \texttt{local-runtime/1} & Changes method domain or result/refusal semantics; run operation conformance\\
Artifact schema & \texttt{causalbank.report/1} and input-role schemas & Changes persisted object meaning; migrate or refuse stale artifacts\\
Wire protocol & MCP \texttt{2025-11-25} or \texttt{2026-07-28} & Changes framing, negotiation, tool result, or session behavior; qualify client/server pair\\
Optional extension & Namespaced negotiated task or other extension & Changes opt-in behavior; require explicit capability agreement\\
\hline
\end{tabular}
\end{center}

The current source pins local request/response schemas at 1.0.0 and reports
operation version \texttt{local-runtime/1}; reports use their own artifact
schema. The MCP adapter is pinned to the 2025-11-25 initialize/initialized
protocol and Python SDK 1.27.0. The pinned 2026-07-28 specification instead
puts version metadata on each request and defines server discovery. Modern
servers must implement discovery; a modern client may call another request
directly and handle a version error. This differs from the recommended probe
for a client that must distinguish modern and legacy servers. The
installed SDK snapshot declares the legacy protocol as its latest version;
modern support is not implemented or trialed here. These are audited source
facts and bounded runtime observations, not a claim that the newer standard is
unsupported by all SDKs.\citep{v12_mcp_legacy_lifecycle,v12_mcp_versioning,v12_mcp_basic,v12_mcp_stdio,exec16_parity_observed}

The protocol transition rule is subtle. A dual-era stdio client probes
\texttt{server/discover}. A discovery result means modern; the client chooses
a mutually supported version. A recognized modern protocol-version error
also means modern: retry with one of its advertised supported versions or
report incompatibility, but do not fall back to legacy initialization. A
different error or timeout identifies the legacy era, so a dual-era client
tries the legacy handshake. The fallback cannot depend on one specific
unknown-method error code because legacy implementations may return several
codes or no response. This is a protocol decision, not a scientific
operation.\citep{v12_mcp_stdio,v12_mcp_versioning}

The authored decision table covers a discovery result with and without a
common version, a recognized unsupported-version error, another recognized
modern error, two ordinary errors, timeout, and a modern-only client. It
retains the no-downgrade rule for every recognized modern error. The table is
not a wire parser and no simulated row establishes actual SDK behavior. Before
advertising dual-era support, a real client/server matrix must test malformed
discovery bodies, empty and contradictory version lists, advertised-version
retry, recognized errors, unknown errors, timeout, cache lifetime, and server
restart. It must also compare the exact canonical response and tool capability
map across each supported protocol.\citep{exec17_continuity_reference,exec17_continuity_contract}

Compatibility must not silently widen the scientific request. A protocol
version negotiates message behavior; it cannot decide whether an interval is
a confidence region, convert a posterior set to a sensitivity region, or
repair an ambiguous treatment. A schema version does not prove assignment
truth, temporal ordering, or graph-class validity. The adapter should reject
unknown mandatory semantic fields, support only explicit migrations, and
retain original request bytes and version identities for audit. Unrecognized
scientific semantics require refusal or a new reviewed operation contract,
not best-effort coercion.
\citep{v12_mcp_basic,exec16_parity_design}

\section{Make adoption evidence exact, external, and conditional}
\label{exec17:signed-adoption}

The evaluation object must identify what was assessed. A proposed release
manifest binds the reader PDF and source audit; canonical schemas and operation
registry; runtime and adapter builds; dependency and protocol versions; the
consumer decoder; the scientific target and sampling law; the declared
platform; and the acceptance criteria. A change to any item starts an impact
review. Prior receipts may be carried forward only when their owner records
which claims remain applicable and why. A hash establishes byte identity, not
who signed those bytes, whether that person had authority, or whether the
scientific statement is true.

The exact authored gate table contains six assumed conditions: independent
scientific review; authorized bank data and truth; actual consumer acceptance;
target-platform acceptance; a matched economic case; and a sponsor decision.
There are $2^6=64$ Boolean combinations, but only the all-true combination is
eligible for scoped adoption under the stated conjunction. This is a logic
check on a proposed rule. It supplies zero authentic receipts, identifies no
signer, and performs no cryptographic verification. The current external-input
register has seven missing groups, including scientific review, bank truth,
platform, consumers, workflow economics, and sponsor decision; the named
owners remain unassigned and the intake ledger contains no external
receipt.\citep{exec17_continuity_reference,exec17_continuity_contract}

Local preparation can still be complete enough to hand over. The package
should contain an exact-release manifest, a criterion-by-criterion evaluation
protocol, a custody-safe bank data manifest, the consumer and platform test
plans, a full-cost ledger, a conflicts statement, and a receipt schema with
review scope, owner role and organization, evidence identity, criteria,
unresolved findings, and assessed release. The intake script can verify
required fields and hashes. It explicitly cannot authenticate a signature,
establish legal access, validate bank truth, or decide scientific adequacy.
Those determinations belong to the authentic responsible owner and reviewer.
The existing external-input contract states this boundary directly.
\citep{exec17_continuity_contract}

The six-field release example changes the reader, core, adapter, consumer,
scientific context, or platform one at a time. Every change breaks exact
release equality and triggers impact review; it does not automatically erase
all prior evidence. A scoped owner may establish that an earlier test remains
applicable, but the evaluator must bind that judgment to the successor
release. Otherwise the stale receipt is not evidence for the changed system.
The proposal deliberately has no success branch that manufactures the missing
owner, signature, bank history, or institutional decision.
\citep{exec17_continuity_reference,exec17_continuity_contract}

The integrated next step is a preserving, generated/public prototype that
compares canonical library behavior, CLI decoding, and MCP tool results under
the same frozen scientific contexts; injects job and storage failures at
declared boundaries; and records full request, response, artifact, cost, and
refusal denominators. If those tests show that semantic preservation and
bounded recovery are useful, request narrowly scoped consumer and platform
trials with authentic owners. Expensive bank experiments follow only after
public evidence shows merit and the bank agrees on a meaningful target and
precision threshold. Independent review and adoption remain separate from
author completion. Until then, the appropriate status is a developed proposal
with specific implementation and institutional acceptance gaps—not a
deployed or approved product.
\label{exec17:integration-continuity-end}
